\documentclass{article}
\usepackage{geometry}
\usepackage{graphicx} % Required for inserting images
\usepackage{amsmath, amsthm, amssymb}
\usepackage{parskip}
\newgeometry{vmargin={15mm}, hmargin={24mm,34mm}}
\theoremstyle{definition} 
\newtheorem{definition}{Definition}

\newtheorem{theorem}{Theorem}[section]
\newtheorem{lemma}[theorem]{Lemma}
\newtheorem{proposition}[theorem]{Proposition}
\newtheorem{corollary}{Corollary}[theorem]


\title{Unique Factorization Domains}
\date{April 2025}
\author{Boran Erol}

\begin{document}

\maketitle

We will show that not every PID is a Euclidean domain by relating the principal
ideal property with a weakening of the Euclidean condition.

\begin{definition}
    Let $ N $ be a \textbf{Dedekind-Hasse norm} if $ N: R \to \mathbb{Z}^{\geq 0} $
    satisfies the following conditions:

    \begin{enumerate}
        \item $ N(a) = 0 $ if and only if $ a = 0_{R} $.
        \item For every nonzero $ a,b \in R $, either $ a \in (b) $ or there's a non-zero
        element $ x \in (a,b) $ such that $ N(x) < N(b) $.
    \end{enumerate}
\end{definition}

The second condition can be rephrased as: either $ a $ divides $ b $ or there exists $ s,t \in R $
such that $ N(as - bt) < N(b) $.

Notice that $ R $ is a Euclidean domain if we can pick $ N $ such that $ s = 1 $ and this condition
is satisfied.

\begin{proposition}
    An integral domain $ R $ is a PID if and only if $ R $ has a Dedekind-Hasse norm.
\end{proposition}
\begin{proof}
    Assume $ R $ has a Dedekind-Hasse norm.

    Let $ I $ be a non-zero ideal in $ R $ and let $ b \in I $ be such that
    $ N(b) $ is minimal. We will show that $ I = (b) $.

    Let $ a \in I $. If $ a = 0 $, we're done.

    Assume that $ 0 \neq a \notin (b) $ by contradiction.
    Then, by the Dedekind-Hasse condition, there exists $ x \in (a,b) $
    such that $ N(x) < N(b) $. But $ (a,b) \subseteq I $, so we have contradicted
    the minimality of $ N(b) $.

    Now, assume that $ R $ is a PID. Then, $ R $ is also a UFD.

    Define $ N $ by letting $ N(0) = 0 $, $ N(u) = 1 $ if $ u $ is a unit,
    and $ N(a) = 2^{n} $ if $ a = p_{1} \ldots p_{n} $ is the irreducible decomposition of $ a $.

    Clearly, $ N $ is multiplicative. Let's now show that $ N $ is a Dedekind-Hasse norm.


\end{proof}



\newpage




\textbf{Let $ R $ be a UFD. Is there a geometric condition that captures this?}

\textbf{Are localizations of $ R $ also UFDs?}



\end{document}
